\section{IV. CONCULUSIONS}

\subsection{Sintesis Temuan}
\begin{frame}{Kesimpulan Utama}
    \begin{itemize}
        \item \textbf{1. Pencapaian:} Implementasi algoritma kuantum efisien untuk reduksi faktor \textit{noisy} HJM.
        \item \textbf{2. Efektivitas:} Solusi untuk masalah finansial kompleks dengan sumber daya terbatas.
        \item \textbf{3. Skalabilitas:} Diklaim memiliki skalabilitas untuk masalah berdimensi lebih tinggi di masa depan.
        \item \textbf{5. Metodologi:} Sebagai pendekatan paling pragmatis dan layak jalan untuk hardware kuantum saat ini.
        \item \textbf{6. Bukti Empiris:} Hasil estimasi \textit{eigenvector} pada matriks $2 \times 2$ yang akurat.
        \item \textbf{8. Pelajaran:} Mengidentifikasi batas kemampuan hardware saat ini sebagai panduan riset masa depan.
        \item \textbf{10. Optimisme:} Menunjukkan bahwa aplikasi praktis komputer kuantum di keuangan sudah di ambang pintu.
        \item \textbf{12. Signifikansi:} Sebagai karya pionir yang menetapkan standar baru dalam bidang \textit{quantum finance}.
    \end{itemize}
\end{frame}

\begin{frame}{Pertanyaan yang dijawab pada Kesimpulan}
    \centering
    \small
    \begin{tabular}{lp{4.5cm}p{6cm}}
        \toprule
        Kalimat ke & Pertanyaan & Jawaban \\
        \midrule
        1 & Pencapaian teknis utama? & Reduksi faktor \textit{noisy} HJM via algoritma efisien. \\
        5 & Justifikasi hibrida? & Pendekatan paling pragmatis untuk hardware saat ini. \\
        8 & Pelajaran 4x4? & Identifikasi batas kemampuan hardware kontemporer. \\
        11 & Dasar optimisme? & Tren perkembangan teknologi hardware yang pesat. \\
        \bottomrule
    \end{tabular}
\end{frame}

\subsection{Transparansi Riset}
\begin{frame}{Open Science}
    \begin{itemize}
        \item \textbf{Repository:} Melalui repositori GitHub publik untuk menjamin transparansi dan replikabilitas.
    \end{itemize}
\end{frame}
